\section*{Bab \ref{ch:b2:two-wave-interference} --- Interferensi Dua Gelombang}

\begin{solution}{exo:b2:two-wave-interference:1}
$i = \lambda D/a = \qty{2.4}{mm}$; $x/i = 4.2$; rumbai gelap kesepuluh di $9.5i = \qty{22.5}{mm}$;
$a = \qty{2}{mm}$: \qty{0.47}{mm}; dalam air $\lambda/n$: \qty{1.8}{mm}.
\end{solution}

\begin{solution}{exo:b2:two-wave-interference:2}
$1.25I_0 \pm I_0$: $V = 0.8$. $1.01I_0 \pm 0.2I_0$: $V = 0.2$ --- suku interferensinya
sejalan dengan nisbah \emph{amplitudo}-nya, $1/10$, bukan nisbah intensitasnya. Tapisnya:
$V = 2\sqrt{0.5}/1.5 = 0.94$.
\end{solution}

\begin{solution}{exo:b2:two-wave-interference:3}
$2ne = \qty{665}{nm}$: terang di $665/(p + \tfrac12)$: \qty{443}{nm} (biru-lembayung);
gelap di $665/p$: \qty{665}{nm} (merah): sebuah selaput biru. $e = \qty{1}{\micro m}$: $2ne = \qty{2660}{nm}$,
terang di $760$, $591$, $484$, \qty{409}{nm} --- empat warna sekaligus, sebuah
putih kemerahjambuan yang tersapu.
\end{solution}

\begin{solution}{exo:b2:two-wave-interference:4}
$r_p = \sqrt{p\lambda R}$: $1.09$, $1.53$, \qty{1.88}{mm}; $p = r^2/\lambda R = 85$ cincin dalam
\qty{1}{cm}; pusatnya gelap: tebalnya nol dan satu pembalikan tanda. Airnya:
jari-jarinya dibagi $\sqrt n$, $98$ cincin.
\end{solution}

\begin{solution}{exo:b2:two-wave-interference:5}
(a) $i = 0.8$ dan \qty{1.4}{mm}; merah kedua \qty{2.8}{mm}, lembayung ketiga \qty{2.4}{mm}.
(b) Orde $p$ merah di $1.4p$ mm berjumpa dengan orde $p + 1$ lembayung di $0.8(p + 1)$: sejak
$p = 2$ ordenya bertumpang tindih. (c) $\delta = ax/D = \qty{2.5}{\micro m}$: yang hilang $\lambda = \delta/(p +
\tfrac12)$: $714$, $556$, \qty{455}{nm}. (d) Pusatnya putih; lalu rumbai dengan
tepi dalam yang biru dan tepi luar yang merah; pada orde ketiga warnanya
bercampur menjadi putih.
\end{solution}

\begin{solution}{exo:b2:two-wave-interference:6}
(a) $b < \lambda L/4a = \qty{59}{\micro m}$. (b) $\pi ab/\lambda L = 6.7$: kontrasnya $|\sin6.7|/6.7
= 0.06$. (c) Sumber selebar sesentimeter merentangi ratusan \omterm{prop:b2:two-wave-interference:young}{jarak antarrumbai}
geseran: tak ada rumbai. (d) $a < \lambda/\theta = \qty{61}{\micro m}$.
\end{solution}

\begin{solution}{exo:b2:two-wave-interference:7}
(a) $a = 2h = \qty{0.4}{mm}$: $i = \lambda D/a = \qty{1.25}{mm}$. (b) Di $x = 0$ lintasannya
sama namun rumbainya gelap: pemantulannya menambahkan $\pi$ (pembalikan
tanda pada datang menyerempet di medium yang lebih rapat). (c) $\delta = ax/D <
\qty{50}{\micro m}$: $100$ rumbai. (d) Lautnya adalah cerminnya, antena di puncak
tebingnya adalah layarnya: ketika sumbernya naik, rumbainya menyapu lewat lalu
memberi ketinggiannya.
\end{solution}

\begin{solution}{exo:b2:two-wave-interference:8}
(a) $2d = 42\lambda$: $d = \qty{12.4}{\micro m}$. (b) $100/42 = \qty{2.4}{mm}$. (c) $2nd = p\lambda$:
$56$ rumbai. (d) Rumbainya beranjak ke ujung jauhnya lalu merenggang ketika $d$
turun; seperempat rumbai adalah $\lambda/8 = \qty{74}{nm}$ pada $d$.
\end{solution}

\begin{solution}{exo:b2:two-wave-interference:9}
(a) Hanya yang pertama (udara $\to$ minyak, lebih rapat); minyak $\to$ air menuju indeks
yang lebih rendah: $\delta = 2ne + \lambda/2$, terang di $2ne = (p + \tfrac12)\lambda$. (b) $2ne =
\qty{1160}{nm}$: yang menguat $773$ (tepi cahaya tampak), \qty{464}{nm} (biru);
yang tertiadakan $1160/p$: \qty{580}{nm} (kuning): biru. (c) $2ne = \lambda/2$: $e = \qty{72}{nm}$.
(d) $\delta = 2ne\cos r$ turun dengan sudutnya: panjang gelombang yang menguat
beranjak ke arah birunya. (e) Kedua pantulannya lalu membalikkan tandanya, tanpa
$\lambda/2$ tambahan: terang di $2ne = p\lambda$, gelap di $(p + \tfrac12)\lambda$ --- yaitu
\omterm{prop:b2:wave-interfaces:optics}{lapisan antipantul} seperempat gelombang.
\end{solution}

\begin{solution}{exo:b2:two-wave-interference:10}
(a) $(n - 1)e = \qty{10}{\micro m}$, $20$ rumbai, ke arah celah yang tertutupi (yang
lintasan geometrinya harus memendek sebesar \qty{10}{\micro m} untuk memulihkan $\delta = 0$):
$x = (n - 1)eD/a = \qty{2}{cm}$. (b) Rumbai putihnya duduk di sana; dispersi
kacanya ($\Delta n \approx 0.01$ di seluruh cahaya tampak, yakni \qty{0.2}{\micro m} lintasan)
berada di bawah $\ell_c$: masih tampak, sedikit berwarna. (c) Ukurlah
geseran rumbai putihnya: $e = (\text{geseran})\,a/(n - 1)D$. (d) $e(n - 1)\theta^2
(1 + 1/n)/2 = \qty{63}{nm}$ pada \qty{5}{\degree}: seperdelapan rumbai.
\end{solution}

\begin{solution}{exo:b2:two-wave-interference:11}
(a) $\theta = \qty{2.3e-7}{rad}$: $a = 1.22\lambda/\theta = \qty{3.1}{m}$. (b) Batasnya $1.22\lambda/d =
\qty{0.057}{''}$ berada di atas diameter bintangnya, dan atmosfernya mengaburkan
sampai $1''$. (c) $1.22\lambda/100 = \qty{7e-9}{rad} = \qty{1.4}{mas}$. (d) Rumbainya menuntut
$\delta < \ell_c = \lambda^2/\Delta\lambda = \qty{32}{\micro m}$.
\end{solution}

\begin{solution}{exo:b2:two-wave-interference:12}
(a) $\delta = 2ne\cos r + \lambda/2$; $\cos r = \sqrt{1 - \sin^2i/n^2} \approx 1 - i^2/2n^2$.
(b) $2ne/\lambda = 10186.8$: pusatnya, dengan $p + \tfrac12$ antara $10186.5$ dan
$10187.5$, tak terang maupun gelap. (c) $i_p^2 = 2n^2[1 - (p + \tfrac12)\lambda/
2ne]$ bagi $p = 10186$, $10185$, $10184$: $i = 0.011$, $0.024$, \qty{0.032}{rad}, $r = fi
= 5.4$, $11.8$, \qty{15.8}{mm}. (d) Sudut $i$ sendirilah yang menetapkan ordenya: setiap
titik sumbernya menyuapi cincin yang sama pada sudut yang sama, sehingga cincinnya, yang terlihat
di tak hingga, bertahan menghadapi sumber yang lebar.
\end{solution}

\begin{solution}{pb:b2:two-wave-interference:1}
\textbf{1.} Tiap cerminnya mencitrakan $S$ lewat kesetangkupan pada jarak $d$ dari
garis sentuhnya; kedua citranya terletak pada lingkaran berjari-jari $d$, pada
sudut $2\alpha$ terpisah: $a = 2\alpha d = \qty{2}{mm}$.

\textbf{2.} $S_1$ dan $S_2$ adalah dua citra getaran yang sama: lompatan
fasenya sama.

\textbf{3.} $i = \lambda(d + D)/a = 589 \times 10^{-9} \times 1.7/2 \times 10^{-3} = \qty{0.50}{mm}$: sepuluh
rumbai di seluruh \qty{5}{mm}.

\textbf{4.} $\delta = ax/(d + D) = \qty{2.9}{\micro m}$ di tepinya, $p = 5$: jauh di bawah
$\ell_c/\lambda = 34\,000$.

\textbf{5.} Geserannya $bD'/d = 20 \times 10^{-6} \times 1.7/0.2 = \qty{0.17}{mm}$, sepertiga
$i$: kontrasnya $|\operatorname{sinc}(\pi ab/\lambda d)| = \operatorname{sinc}(1.07) = 0.82$.

\textbf{6.} $\pi ab/\lambda d = 5.3$: $|\sin5.3|/5.3 = 0.16$.

\textbf{7.} Amplitudonya $\sqrt{0.9}$, $\sqrt{0.8}$: $V = 2\sqrt{0.72}/1.7 = 0.998$.

\textbf{8.} $(n - 1)e = \qty{5}{\micro m}$: $8.5$ rumbai; sepele diikuti dengan
cahaya natrium; dalam cahaya putih rumbai putihnya beranjak sebesar $8.5i = \qty{4.2}{mm}$
lalu masih ada di sana (dispersi kacanya memakan sepersekian
mikrometer).

\textbf{9.} Sebuah rumbai pusat yang putih, dua atau tiga rumbai berwarna di
tiap sisinya, lalu sebuah putih yang seragam.

\textbf{10.} Udara $\to$ minyak membalikkan, minyak $\to$ air (indeks lebih rendah) tidak:
$\delta = 2ne + \lambda/2$; terang bagi $2ne = (p + \tfrac12)\lambda$.

\textbf{11.} $e = \lambda/4n$: merah \qty{112}{nm}, hijau \qty{95}{nm}, lembayung \qty{72}{nm}.

\textbf{12.} $2ne = \qty{1740}{nm}$: yang menguat $1740/(p + \tfrac12)$: $696$ (merah), $497$
(biru-hijau), \qty{387}{nm}; yang tertiadakan $1740/p$: $580$ (kuning), \qty{435}{nm}:
sebuah campuran keunguan merah dan biru-hijau.

\textbf{13.} Tebalnya berubah sepanjang jalannya; tiap warna terang
di tempat $2ne$ bernilai tepat, sehingga pitanya merupakan kontur ketebalan
yang sama.

\textbf{14.} Ordenya turun: warnanya menjelajahi runtunannya
ke arah ujung selaput tipisnya; pada $e \to 0$ pembalikan tanda tunggalnya menyisakan
$\delta = \lambda/2$ bagi setiap warna: gelap --- seperti bagi selaput sabun dalam udara.

\textbf{15.} Beberapa nanometer: $2ne \ll \lambda$ bagi semua warnanya, kedua
pantulannya berlawanan fase lalu meniadakan: hitam.

\textbf{16.} $\sin r = \sin60^\circ/1.45$, $r = \qty{37}{\degree}$, $\cos r = 0.80$: $2ne\cos r
= \qty{1390}{nm}$, terang di $928$, $557$, \qty{398}{nm}: merah pertanyaan 12 telah
berubah hijau-kuning --- ke arah birunya.

\textbf{17.} Ketika $2ne$ melampaui $\ell_c \approx \qty{1}{\micro m}$, yakni di luar sekitar
\qty{0.4}{\micro m} minyak, ordenya bertumpang tindih dan warnanya memudar menjadi kelabu.

\textbf{18.} Kedua sinar yang terpantul dari selaputnya bagi arah sumber
mana pun berjumpa kembali pada selaputnya: rumbainya terlokalkan di sana dan
sumber yang lebar dan berarah banyak sekadar menambahkan pola yang serupa.

\textbf{19.} $\ell_c/\lambda$: $34\,000$ dengan natrium, $1.7 \times 10^7$ dengan lasernya.

\textbf{20.} $b < \lambda d/4a$: \qty{15}{\micro m} bagi cerminnya, \qty{59}{\micro m} bagi
celahnya.

\textbf{21.} Gambarlah dua titik sumber: bagi masing-masingnya, kedua sinar pantulnya
bergabung kembali di titik yang sama pada selaputnya, dengan $\delta = 2ne\cos r$ yang sama
pada datang hampir tegak lurus; pola selaputnya tak beranjak, pola Young
beranjak.

\textbf{22.} $\lambda/\Delta\lambda = 589/0.02 \approx 30\,000$.

\textbf{23.} $550/30 \approx 18$ rumbai.

\textbf{24.} Cahaya koheren yang dihamburkan butir layarnya berinterferensi
dengan fase acak di matanya: bercak.

\textbf{25.} Cerminnya: \omterm{prop:b2:two-wave-interference:young}{pembelahan muka gelombang}, rumbai di mana pun
(tak terlokalkan), jumlahnya dibatasi $\ell_c/\lambda$, sumber yang lebar menghapusnya.
Selaputnya: \omterm{prop:b2:two-wave-interference:film}{pembelahan amplitudo}, rumbai pada selaputnya, warnanya
dibatasi $\ell_c \sim \qty{1}{\micro m}$, sumber yang lebar tak berbahaya.
\end{solution}
